% Independently reviewed by root and cohomology; ready for integration.
% E203--E204 reconstruct finite inputs; they do not certify A4.57's bound.
\subsection{The exact compact quaternion detector reconstructed from words}

\begin{result}{E.203}{Proved}
{The compact detector has the stated order and exact coupled CM residues}
We independently reconstruct the order and CM-centralizer clauses of
historical GT A.4.18.  Let $\Sigma=F(a,b)$, put
$c=[a,b]=aba^{-1}b^{-1}$, and define
\[
 \begin{gathered}
 q_8:\Sigma\twoheadrightarrow Q_8,\qquad a\mapsto i,\quad b\mapsto j,
 \qquad N_8=\ker q_8,\\
 T=N_8\cap\ker\bigl(\Sigma\to(\mathbb Z/3)^2\bigr),\qquad
 L_8=\gamma_3(N_8)N_8^3,\\
 L_c=L_8T'\langle\!\langle c^2\rangle\!\rangle_\Sigma,
 \qquad R_c=L_c\cap\ker\bigl(\Sigma\to(\mathbb Z/12)^2\bigr),\\
 P_c=\Sigma/R_c.
 \end{gathered}                                                   \tag{E.203.1}
\]
All subgroups here are literal subgroups of the indicated free group;
their closures give the same finite quotients in its full profinite
completion.  Then
\[
 |\Sigma/L_c|=8\cdot3^{14},\qquad
 |P_c|=32\cdot3^{14}=153055008,
 \qquad P_c^{\mathrm{ab}}=(\mathbb Z/12)^2.                       \tag{E.203.2}
\]
Consider the ordinary surface automorphisms
\[
 \begin{array}{c|cc|c}
       &a&b&\text{period}\\\hline
 u&a&ba&U=\left(\begin{smallmatrix}1&1\\0&1\end{smallmatrix}\right)\\
 l&ab^{-1}&b&L=\left(\begin{smallmatrix}1&0\\-1&1\end{smallmatrix}\right).
 \end{array}                                                     \tag{E.203.3}
\]
Their action on the characteristic three-subgroup of $P_c$, joined
to abelian reduction modulo twelve and divided by the actual inner
surface actions, has an explicitly reconstructed image $J$ of order
\[
                     |J|=1152\cdot3^4=93312.                     \tag{E.203.4}
\]
Write $r$ for the image of $ul$ and $R=UL$.  The period image of
$C_J(r)$ consists exactly of
\[
 q\bmod4=R^i\bmod4,\qquad q\bmod3=R^j\bmod3,
 \qquad i,j\in\mathbb Z/6,\qquad i\equiv j\pmod2.                \tag{E.203.5}
\]
Each of these eighteen residues has nine centralizing lifts.  Thus
$|C_J(r)|=162$.  In particular the residue $5I$, whose two exponents
are $0$ and $3$, is excluded.  This is a computation on the actual
surface-word image, not an enumeration of formal unrelated matrices.
\end{result}

\begin{proof}
We give the complete finite model and the exhaustiveness argument
for the certificate.  No historical matrix file is imported.

There are twenty-four ordered generating pairs in $Q_8$.  The
assignment $(i,j)$ to any such pair extends uniquely to an
automorphism: the two images have square $-1$ and anticommute.
Consequently every epimorphism $\Sigma\to Q_8$ has the same kernel.
This proves characteristicness of $N_8$.  The subgroup $T$ is
characteristic as well.  Both automorphisms in (E.203.3) fix $c$
as an element.  They therefore preserve all the groups in
(E.203.1), including the final abelian join.  These are the usual
generators of the orientation-preserving ordinary outer action:
the exact discrete identification
$\operatorname{Out}(F_2)=\operatorname{GL}_2(\mathbb Z)$ is recalled
in \citet[\S3, p.335]{LitBuxErshovRapinchuk2011}.

\paragraph{The literal free-cover model.}
Breadth-first search in the eight-element quaternion Cayley graph,
with letter order $a,b,a^{-1},b^{-1}$, gives a free Schreier basis
of $N_8$ with nine words.  In that deterministic order they are
\[
 \begin{gathered}
 a^4,\quad a^2b^2,\quad a^{-1}b^2a^{-1},\quad
 baba^{-1},\quad b^2a^{-2},\\
 b^{-1}ab^{-1}a^{-1},\quad abab^{-1},\quad
 ab^2a,\quad ab^{-1}ab.
 \end{gathered}                                                   \tag{E.203.6}
\]
The corresponding graph has sixteen positive edges and eight
vertices, so the rank is nine.  Rewrite any word in $N_8$ by
following its actual quaternion vertices.  This gives its exact
word in these nine generators, including negative edges.

Let $V_9=\mathbb F_3^9$.  The free class-two exponent-three group on
these generators is the set
\[
 V_9\oplus\bigwedge^2V_9,\qquad
 (v,z)(w,t)=(v+w,z+t+2v\wedge w).                                 \tag{E.203.7}
\]
Here $2=1/2$ in $\mathbb F_3$.  The commutator is $(0,v\wedge w)$,
and inversion is negation.  Collecting a word proves the universal
property: every class-two exponent-three group generated by nine
elements is an image of (E.203.7).  Conversely (E.203.7) has the
stated $3^{45}$ elements and satisfies those relations.  It is
therefore exactly $N_8/L_8$, rather than merely an approximation to
that quotient.

The two ordinary exponent sums, reduced modulo three, give a
surjective map $V_9\to\mathbb F_3^2$.  If $V_T$ is its
seven-dimensional kernel, the image of $T'$ in (E.203.7) is
exactly $\bigwedge^2V_T$, of dimension twenty-one.  Indeed every
vector of $V_T$ is represented by a word in $T$, and the
commutator in (E.203.7) depends only on the two abelian vectors.
No central-coordinate restriction on those representatives is
being imposed.

We next rewrite the actual word $c^2$ and take the smallest
subspace containing its logarithm and $\bigwedge^2V_T$, stable
under the two actual deck substitutions
$\operatorname{Ad}(a),\operatorname{Ad}(b)$ and under brackets
with the nine generators.  This is exactly its normal closure in
the class-two group.  Subgroup addition and Lie addition agree
after retaining brackets in odd characteristic; the two deck
substitutions and conjugations inside $N_8$ generate all
conjugations by $\Sigma$.  Thus there is no supplementary
normality condition.  The stable relation space has dimension
thirty-one.  Its quotient has dimension fourteen, with
six-dimensional abelianization and eight-dimensional derived
space.  The program retains the complete $31\times45$ relation
matrix and the $14\times45$ quotient map.  It checks that all
four substitutions $u,l,\operatorname{Ad}(a),\operatorname{Ad}(b)$
preserve that relation space.  This proves
$|N_8/L_c|=3^{14}$.

The exponent lattice of $N_8$ is exactly $2\mathbb Z^2$.
Containment follows from $Q_8^{\mathrm{ab}}=(\mathbb Z/2)^2$;
the words $a^2c^{-1}$ and $b^2c^{-1}$ prove the reverse
containment.  Hence the exponent lattice of $L_8$ is
$6\mathbb Z^2$.  The additional generators of $T'$ and of the
peripheral normal closure have zero exponent sums, so $L_c$ has
that same lattice.  Joining reduction modulo twelve adds
exactly four.  The exponent lattice of $R_c$ is
$12\mathbb Z^2$.  This proves all assertions in (E.203.2).

\paragraph{The complete image computation.}
Choose six quotient coordinates complementary to the derived
space, followed by eight derived coordinates.  The four actual
substitutions become matrices
\[
                  \begin{pmatrix}A&0\\ C&B\end{pmatrix},
          \qquad A\in\operatorname{GL}_6(3),\quad
                 B\in\operatorname{GL}_8(3).                    \tag{E.203.8}
\]
Their entries are obtained by substituting the literal words
(E.203.6) into (E.203.3) and the two inner substitutions, then
applying the quotient map just constructed.  The eight central
coordinates are generated by brackets, so $B$ is determined by
$A$.  The space of inner derivations has rank six; in particular
the centre of this fourteen-dimensional Lie algebra is exactly
its eight-dimensional derived space.

An automorphism with $A=B=1$ is an additive element of
$\operatorname{Hom}(\mathbb F_3^6,\mathbb F_3^8)$.  Quotienting
by the actual inner derivations gives a forty-two-dimensional
vector space.  The complete finite linear base, retaining $A$
and the signed period modulo twelve, has $9216$ elements.
Breadth-first multiplication by the four actual substitutions
lists every element and retains a source-word representative.
All $9216\cdot4=36864$ Cayley--Schreier discrepancies are then
computed in that forty-two-dimensional space.  Their span has
dimension four and is invariant under the four source actions.
Schreier generation proves that this is the entire kernel, not
just a lower bound: a word with trivial base is a product of
these discrepancies, and the discrepancy group is additive.

The inner deck actions, after dividing by inner derivations,
have order eight.  Their action on the six-dimensional abelian
space is faithful; the certificate lists all eight distinct
matrices.  They are a normal subgroup of the joint image,
because every source automorphism conjugates an inner surface
automorphism to the corresponding inner automorphism.  Dividing
the complete image by this subgroup gives
$9216\cdot3^4/8=93312$, as asserted.

\paragraph{The complete centralizer computation.}
The period image consists of all $1152$ elements of
$\operatorname{SL}_2(\mathbb Z/12)$.  Exactly thirty-six
commute with $R$.  They have the unique descriptions in
(E.203.5) before the parity condition is imposed.  For each,
choose its retained source-word representative and multiply it
by every one of the $3^4$ elements of the computed kernel.
Form its commutator with the actual matrix of $ul$.  Its linear
base, when centralization is possible, is one of the eight
actual deck bases.  Remove that named deck representative and
test whether the remaining central-coordinate matrix is an
inner derivation.  This is exactly equality to one in $J$.

The resulting numbers of centralizing lifts are the following
complete table; rows are $i$ and columns are $j$:
\[
 \begin{array}{c|rrrrrr}
     &0&1&2&3&4&5\\\hline
 0&9&0&9&0&9&0\\
 1&0&9&0&9&0&9\\
 2&9&0&9&0&9&0\\
 3&0&9&0&9&0&9\\
 4&9&0&9&0&9&0\\
 5&0&9&0&9&0&9
 \end{array}                                                     \tag{E.203.9}
\]
The code checks every residue and every kernel vector, retaining
all source words and the four-dimensional kernel basis.  Thus
there is no untested lift or omitted ordinary coset.

\paragraph{The reproducible certificate.}
The independent program
\path{code/check_compact_quaternion_detector.py} contains the
quaternion multiplication, free-word rewriting, class-two law,
exact row reduction over $\mathbb F_3$, quotient maps, actual
automorphism substitutions, complete base enumeration and all
centralizer tests.  It uses no floating-point arithmetic and
imports no historical certificate.  Its JSON retains the basis,
all four $14\times14$ matrices, full relation and quotient
matrices, inner-derivation basis, kernel basis, and the thirty-six
source representatives.  The final payload SHA256 is
\[
 \begin{gathered}
 \texttt{868ede05a7e14876221494203a607efef}\\[-2pt]
 \texttt{6a1d3e982d93207feeac268bdac5d6f}.
 \end{gathered}                                                   \tag{E.203.10}
\]
This provides the finite inputs to the preceding exact group
arguments.  It does not assert the other homology-module
identifications in the historical A.4.18, or its later compiler
bounds, without their own proofs.
\end{proof}

\begin{result}{E.204}{Proved}
{The detector is the actual ordinary outer image and retains the ordered-pole module}
For the group $P_c$ of E.203, the ordinary modular image in
$\operatorname{Out}(P_c)$ is isomorphic to the computed group
$J$.  Thus the CM-centralizer conclusion of E.203 holds in the
actual ordinary outer image itself.  This proves the
ordinary-image conclusion needed from historical A.4.22 by a
direct coprime argument, without its missing hermitian
normalizer certificate.

Let $\Pi_*=F(A,B,C)$ be the actual two-punctured torus group
with the ordered-pole marking below.  The historical joint
quotient
\[
 W=\operatorname{im}\left(
 \Pi_*\longrightarrow(\mathbb Z/12)^3\times P_c\right),
 \qquad g\longmapsto(g^{\mathrm{ab}},g|_{C=1}),
                                                                    \tag{E.204.1}
\]
is exactly $P_c\times\mathbb Z/12$.  Its actual ordinary
two-punctured-torus outer image is
\[
        J_W=(\mathbb Z/12)^2\rtimes J,\qquad
        |J_W|=144\cdot93312=13436928.                              \tag{E.204.2}
\]
The action on the first factor is the signed compact period
modulo twelve.  The quotient $J_W\to J$ and the full primitive
point-pushing module are retained together, not independently
chosen.  These conclusions verify the finite target and
linearization inputs used in FCP.L1--L3.  They do not establish
the missing bounded branch-covariance assertion BCV.8--10.
\end{result}

\begin{proof}
Put $N=N_8/L_c$.  The kernel $K$ of
$P_c\to\Sigma/L_c$ has order four and is central: for
$s\in\Sigma$ and $t\in L_c$, the commutator $[s,t]$ lies
in $L_c\cap\Sigma'\subseteq R_c$.  Its map to
$P_c^{\mathrm{ab}}=(\mathbb Z/12)^2$ is injective, with
image $6\mathbb Z^2/12\mathbb Z^2$.

The inverse image of $N$ in $P_c$ is a central extension of
the three-group $N$ by $K$.  A Sylow three-subgroup projects
isomorphically to $N$; multiplying it by the central $K$
gives the entire inverse image.  It is consequently normal
and is the unique Sylow three-subgroup there.  Since the
inverse image is normal in $P_c$, this three-subgroup is
normal in $P_c$ as well.  We identify it with $N$.
It is $O_3(P_c)$, hence characteristic.

The faithful quaternion deck action on $N^{\mathrm{ab}}$
verified in E.203 implies
\[
                         C_{P_c}(N)=K\times Z(N).                \tag{E.204.3}
\]
Indeed a centralizing element must have trivial quaternion
label, since inner $N$ actions are trivial on
$N^{\mathrm{ab}}$.  It therefore lies in $K\times N$,
and its $N$ coordinate must be central.  The reverse
inclusion is immediate.

Now take an ordinary automorphism whose class maps to one
in the detected group $J$.  Its period modulo twelve is
the identity, and its restriction to $N$ is the
restriction of an actual inner $P_c$ action.  Remove
that inner action.  The resulting automorphism $\phi$
fixes $N$ pointwise and acts trivially on $P_c^{\mathrm{ab}}$.
For $g\in P_c$,
$\delta(g)=\phi(g)g^{-1}$ belongs to $C_{P_c}(N)$.
Its $K$ component is trivial, because $K$ injects into
the two-primary part of the abelianization while $Z(N)$
has three-power order.  Thus $\delta(g)\in Z(N)$.
It factors through $P_c/N$, a group of order thirty-two,
and is an additive crossed homomorphism for that group's
action on $Z(N)$.

For completeness its triviality in cohomology follows by
literal averaging.  If
$s=\sum_{h\in P_c/N}\delta(h)$, then the cocycle equation
gives
\[
          32\delta(g)=s-gs.
\]
Multiplication by thirty-two is invertible on the finite
three-group $Z(N)$, so $\delta(g)=z-gz$ with $z=32^{-1}s$.
Hence $\phi=\operatorname{Ad}(z)$.  The detected ordinary
outer map has trivial kernel.  Its surjectivity is its
definition, proving the first assertion.

\paragraph{The literal ordered-pole marking.}
Let
$\Omega=C_2(t_1)*C_2(t_2)*C_2(t_3)*C_2(t_4)$ and let
$\Pi_*$ be the kernel of the character taking each
$t_i$ to one.  Its rank is three, and
\[
 \begin{gathered}
 u=t_2t_1,\qquad v=t_3t_1,\qquad z=t_4t_1,\\
 A=u,\qquad B=u^{-1}v,
 \qquad C=u^{-1}vuv^{-1}zv^{-1}u,\\
 u=A,\qquad v=AB,\qquad z=[A,B]CB
 \end{gathered}                                                   \tag{E.204.4}
\]
are mutually inverse free bases.  With $t=t_1t_2t_3t_4$,
the two boundary words are
$t=([A,B]C)^{-1}$ and $C=B(t_1tt_1)B^{-1}$.
These fix the two pole orientations.  The three positive
Artin generators send
$(t_i,t_{i+1})$ to $(t_it_{i+1}t_i,t_i)$.
Rewriting their actual words in (E.204.4) gives
\[
 \begin{array}{c|ccc}
 &A&B&C\\\hline
 a&A&BA&C\\
 b&AB^{-1}&B&C\\
 c&A&C^{-1}BA&hCh^{-1}
 \end{array}
 \qquad h=C^{-1}BAB^{-1}.                                        \tag{E.204.5}
\]
The free-word certificate checks these formulas, both
adjacent braid relations, $ac=ca$, and the full twist
as $\operatorname{Ad}(t)$.

Because $P_c^{\mathrm{ab}}=(\mathbb Z/12)^2$, the first
two abelian coordinates in (E.204.1) are already
determined by its $P_c$ coordinate.  The remaining
coordinate is the exponent of $C$.  The images of
$A,B,C$ generate $P_c\times\mathbb Z/12$, which proves
the claimed exact isomorphism and gives
$|W|=12\cdot32\cdot3^{14}=1836660096$.

In this direct product $C$ is central.  Equation (E.204.5)
therefore shows that every generated automorphism has
the form
\[
        (p,z)\longmapsto
          \bigl(\psi(p),z+\ell(p^{\mathrm{ab}})\bigr),
 \qquad \ell\in\operatorname{Hom}((\mathbb Z/12)^2,\mathbb Z/12).
                                                                    \tag{E.204.6}
\]
An inner automorphism of $W$ is exactly an inner action
on $P_c$ with zero row $\ell$.  The subgroup generated
by $a,b$ has row zero and gives precisely the ordinary
image $J$ already identified.  The actual point-pushing
words
\[
                \xi_1=ca^{-1},\qquad
                \xi_2=aba c^{-1}b^{-1}a^{-1}
\]
have identity compact blocks and rows $(0,-1)$ and
$(1,0)$.  The certificate verifies this in the integral
three-dimensional homology representation before any
finite reduction.  They consequently generate all
$144$ rows, and a nonzero row cannot be inner.  This
proves both exactness and the splitting in (E.204.2).
The full four-strand twist is inner by (E.204.5), so
this is the ordinary $\Gamma_{1,2}$ outer image with
the ordered-pole convention used in E.183 and E.190.

Finally put
$J_0=\left(\begin{smallmatrix}0&-1\\1&0\end{smallmatrix}\right)$
and identify a column $v$ with the row $v^{\mathsf T}J_0$.
Conjugation by a compact period $P$ changes the row to
$v^{\mathsf T}J_0P^{-1}=(Pv)^{\mathsf T}J_0$, since
$P^{\mathsf T}J_0P=J_0$.  This proves the exact signed
action in the semidirect product.  It is an ordinary
finite image construction; no assertion that a general
transported gentle action preserves $\ker(\Pi_*\to W)$
has entered it.
\end{proof}
